\documentclass[11pt,a4paper]{article}
\usepackage{subcaption}

% ============================================================
% PACOTES
% ============================================================
\usepackage[utf8]{inputenc}
\usepackage[T1]{fontenc}
\usepackage[brazilian]{babel}
\usepackage{pgfplots}
\pgfplotsset{compat=1.18}

\usepackage{amsmath,amssymb,amsfonts,mathtools}
\usepackage{bm}
\usepackage{physics}
\usepackage{siunitx}
\usepackage{booktabs}
\usepackage{geometry}
\usepackage{float}
\usepackage{hyperref}
\usepackage{tikz}

\geometry{margin=2.5cm}

\hypersetup{
    colorlinks=true,
    linkcolor=blue,
    urlcolor=blue
}

% ============================================================
% COMANDOS
% ============================================================
\newcommand{\pdvT}[2]{\left(\frac{\partial #1}{\partial #2}\right)_T}
\newcommand{\pdvP}[2]{\left(\frac{\partial #1}{\partial #2}\right)_p}
\newcommand{\pd}[2]{\frac{\partial #1}{\partial #2}}
\newcommand{\pdd}[2]{\frac{\partial^2 #1}{\partial #2^2}}
\newcommand{\e}[1]{\mathrm{e}^{#1}}
\newcommand{\avg}[1]{\left\langle #1\right\rangle}

% ============================================================
% DOCUMENTO
% ============================================================
\begin{document}

\begin{center}
    {\Large \textbf{Quarta Série de Exercícios}}\\[0.6cm]
    {\large Tópicos de Mecânica Estatística -- Transições de Fases -- 2026}\\[0.2cm]
    {\large Modelo de Ising, dualidade e teoria de escala -- 19/09/2026}\\[0.4cm]
    {\large Data de Entrega -- 28/09/2026}\\[0.2cm]
    {\Large \textbf{Autor:} Thiago Siqueira Domingues (Nº USP: 8942194)}
\end{center}

\vspace{0.3cm}

% ============================================================
\section*{Problema 1 -- Forma de escala da energia livre de Gibbs}
% ============================================================

Considere a parte singular da energia livre de Gibbs por partícula de um
fluido, ligeiramente acima da temperatura crítica, na forma de escala
\begin{equation}
    g_s(t,\pi)
    \sim
    t^{2-\alpha}
    F\!\left(\frac{\pi}{t^{\Delta}}\right),
    \label{eq:g_scaling}
\end{equation}
onde
\begin{equation}
    t = \frac{T-T_c}{T_c},
    \qquad
    \pi = \frac{p-p_c}{p_c},
    \qquad
    t \to 0^+,
    \label{eq:variables}
\end{equation}
e \(F\) é uma função bem comportada. Vamos determinar, no diagrama de fases $p-T$, ao longo do caminho crítico \(p=p_c\) (portanto, \(\pi=0\)), os comportamentos críticos das grandezas termodinâmicas:
%
\begin{enumerate}
    \item entropia;
    \item compressibilidade isotérmica;
    \item calor específico a volume constante;
    \item calor específico a pressão constante;
    \item coeficiente de dilatação térmica.
\end{enumerate}
%
Como estamos trabalhando com a energia livre de Gibbs, usamos as relações
termodinâmicas
%
\begin{align}
    s &= -\left(\frac{\partial g}{\partial T}\right)_p,
    \label{eq:s}\\
    v &= \left(\frac{\partial g}{\partial p}\right)_T,
    \label{eq:v}\\
    \kappa_T
    &= -\frac{1}{v}
    \left(\frac{\partial v}{\partial p}\right)_T,
    \label{eq:kappa}\\
    c_p
    &= T\left(\frac{\partial s}{\partial T}\right)_p,
    \label{eq:cp}\\
    \alpha_T
    &= \frac{1}{v}
    \left(\frac{\partial v}{\partial T}\right)_p.
    \label{eq:alphaT}
\end{align}
%
Aqui usamos \(\alpha_T\) para o coeficiente de dilatação térmica, de modo a não
confundi-lo com o expoente crítico \(\alpha\) que aparece em \eqref{eq:g_scaling}. Agora, definimos a variável de escala
%
\begin{equation}
    x \equiv \frac{\pi}{t^{\Delta}},
    \label{eq:x}
\end{equation}
%
de modo que
%
\begin{equation}
    g_s(t,\pi) \sim t^{2-\alpha}F(x).
    \label{eq:gx}
\end{equation}
%
Como
%
\begin{equation}
    \frac{\partial t}{\partial T}=\frac{1}{T_c},
    \qquad
    \frac{\partial \pi}{\partial p}=\frac{1}{p_c},
    \label{eq:dt_dT_dpi_dp}
\end{equation}
%
precisamos calcular as derivadas de \(g_s\) em relação a \(T\) e \(p\). Primeiro,
%
\begin{align}
    \pd{g_s}{t}
    &\sim
    \pd{}{t}\left[t^{2-\alpha}F(x)\right] \nonumber\\
    &\sim
    (2-\alpha)t^{1-\alpha}F(x)
    + t^{2-\alpha}F'(x)\pd{x}{t}.
\end{align}
%
Como
%
\begin{equation}
    x=\pi t^{-\Delta}
    \quad\Longrightarrow\quad
    \left(\frac{\partial x}{\partial t}\right)_\pi
    =-\Delta\frac{x}{t},
\end{equation}
%
obtemos
%
\begin{equation}
    \pd{g_s}{t}
    \sim
    t^{1-\alpha}
    \left[(2-\alpha)F(x)-\Delta xF'(x)\right].
    \label{eq:dgdt_general}
\end{equation}
%
Ao longo de \(p=p_c\), temos \(\pi=0\) e, portanto, \(x=0\). Definindo
%
\begin{equation}
    F_0\equiv F(0),
    \qquad
    F_1\equiv F'(0),
    \qquad
    F_2\equiv F''(0),
    \label{eq:Fcoeff}
\end{equation}
%
segue que
%
\begin{equation}
    g_s(t,0)\sim F_0 t^{2-\alpha}
\end{equation}
%
e
%
\begin{equation}
    \left.\pd{g_s}{t}\right|_{\pi=0}
    \sim
    (2-\alpha)F_0 t^{1-\alpha}.
    \label{eq:dgdt_pc}
\end{equation}
%
Para as derivadas em \(p\), devemos lembrar que
%
\begin{equation}
    \frac{\partial x}{\partial \pi}=t^{-\Delta}.
\end{equation}
%
Para a primeira derivada em \(p\), aplicando a regra da cadeia,
%
\begin{align}
    \pd{g_s}{p}
    &=
    \pd{}{p}\left[t^{2-\alpha}F(x)\right] \nonumber\\
    &=
    t^{2-\alpha}F'(x)\,\pd{x}{p} \nonumber\\
    &=
    t^{2-\alpha}F'(x)\,
    \pd{x}{\pi}\,\pd{\pi}{p} \nonumber\\
    &=
    t^{2-\alpha}F'(x)\,
    t^{-\Delta}\,\frac{1}{p_c} \nonumber\\
    &=
    \frac{1}{p_c}\,
    t^{\,2-\alpha-\Delta}\,
    F'(x).
    \label{eq:dgp_general}
\end{align}
%
Com a notação \(F_1\equiv F'(0)\), segue
%
\begin{equation}
    \left.\pd{g_s}{p}\right|_{p=p_c}
    \sim
    \frac{F_1}{p_c}\,t^{\,2-\alpha-\Delta}.
    \label{eq:dgp}
\end{equation}
%
Para a segunda derivada em \(p\), derivando \eqref{eq:dgp_general} mais uma vez em relação a \(p\),
%
\begin{align}
    \pdd{g_s}{p}
    &=
    \pd{}{p}\left[
    \frac{1}{p_c}\,
    t^{\,2-\alpha-\Delta}\,
    F'(x)
    \right] \nonumber\\
    &=
    \frac{1}{p_c}\,
    t^{\,2-\alpha-\Delta}\,
    F''(x)\,
    \pd{x}{p} \nonumber\\
    &=
    \frac{1}{p_c}\,
    t^{\,2-\alpha-\Delta}\,
    F''(x)\,
    \pd{x}{\pi}\,\pd{\pi}{p} \nonumber\\
    &=
    \frac{1}{p_c}\,
    t^{\,2-\alpha-\Delta}\,
    F''(x)\,
    t^{-\Delta}\,
    \frac{1}{p_c} \nonumber\\
    &=
    \frac{1}{p_c^{2}}\,
    t^{\,2-\alpha-2\Delta}\,
    F''(x).
    \label{eq:d2gp_general}
\end{align}
%
Ao longo de \(p=p_c\) (\(x=0\)), com \(F_2\equiv F''(0)\),
%
\begin{equation}
    \left.\pdd{g_s}{p}\right|_{p=p_c}
    \sim
    \frac{F_2}{p_c^{2}}\,t^{\,2-\alpha-2\Delta}.
    \label{eq:d2gp}
\end{equation}
%
Agora, para a derivada mista \(\partial^{2} g_s/\partial p\,\partial T\), partimos da expressão geral \eqref{eq:dgp_general} e derivamos em relação a \(T\) (com \(p\), e portanto \(\pi\), fixos), e usando
%
\begin{equation}
    \frac{\partial t}{\partial T}=\frac{1}{T_c},
    \qquad
    \left(\frac{\partial x}{\partial T}\right)_{\!\pi}
    =
    \pi\,\frac{\partial}{\partial T}\left(t^{-\Delta}\right)
    =
    -\Delta\,\pi\,t^{-\Delta-1}\,\frac{1}{T_c}
    =
    -\frac{\Delta}{T_c}\,\frac{x}{t},
\end{equation}
%
obtemos
%
\begin{align}
    \frac{\partial^{2} g_s}{\partial p\,\partial T}
    &=
    \frac{\partial^{2} g_s}{\partial T\partial p}
    =
    \frac{1}{p_c}
    \frac{\partial}{\partial T}
    \left[
    t^{\,2-\alpha-\Delta}\,
    F'(x)
    \right] \nonumber\\
    &=
    \frac{1}{p_c}
    \left[
    \frac{\partial}{\partial T}
    \left(t^{\,2-\alpha-\Delta}\right)
    F'(x)
    +
    t^{\,2-\alpha-\Delta}\,
    \frac{\partial F'(x)}{\partial T}
    \right] \nonumber\\
    &=
    \frac{1}{p_c}
    \left[
    \frac{2-\alpha-\Delta}{T_c}\,
    t^{\,1-\alpha-\Delta}\,
    F'(x)
    +
    t^{\,2-\alpha-\Delta}\,
    F''(x)\,
    \left(-\frac{\Delta}{T_c}\,\frac{x}{t}\right)
    \right] \nonumber\\
    &=
    \frac{t^{\,1-\alpha-\Delta}}{p_c\,T_c}
    \left[
    (2-\alpha-\Delta)\,F'(x)
    -
    \Delta\,x\,F''(x)
    \right].
    \label{eq:d2gpt_geral}
\end{align}
%
Ao longo de \(p=p_c\) (\(x=0\)), o termo proporcional a \(x F''(x)\)
desaparece e resta
%
\begin{equation}
    \left.\frac{\partial^{2} g_s}{\partial p\,\partial T}\right|_{p=p_c}
    \sim
    \frac{2-\alpha-\Delta}{p_c\,T_c}\,
    F_1\,t^{\,1-\alpha-\Delta}.
    \label{eq:d2gpt}
\end{equation}
%
\subsection*{(i) Expoente crítico associado a entropia}
%
Da definição \(s=-(\partial g/\partial T)_p\), a parte singular da entropia é
%
\begin{equation}
    s_s
    =-\left(\frac{\partial g_s}{\partial T}\right)_p
    =-\frac{1}{T_c}\pd{g_s}{t}.
\end{equation}
%
Usando \eqref{eq:dgdt_pc}, encontramos
%
\begin{equation}
    s_s(p=p_c)
    \sim
    -\frac{2-\alpha}{T_c}F_0\,t^{1-\alpha}.
    \label{eq:entropy_scaling}
\end{equation}
%
Portanto,
%
\begin{equation}
    \boxed{s_s\sim t^{1-\alpha}}
\end{equation}
%
\subsection*{(ii) Expoente crítico associado a compressibilidade isotérmica}
% ------------------------------------------------------------
A compressibilidade isotérmica é
%
\begin{equation}
    \kappa_T
    =-\frac{1}{v}
    \left(\frac{\partial v}{\partial p}\right)_T
    =-\frac{1}{v}\left(\frac{\partial^2 g}{\partial p^2}\right)_T.
\end{equation}
%
usando \eqref{eq:d2gp} com \(v\to v_c\) no prefator,
%
\begin{equation}
    \kappa_{T,s}
    \sim
    -\frac{F_2}{v_c p_c^2}
    t^{2-\alpha-2\Delta}.
\end{equation}
%
Definindo o expoente crítico \(\gamma\) por \(\kappa_T\sim t^{-\gamma}\),
comparando as potências de \(t\),
%
\begin{equation}
    -\gamma=2-\alpha-2\Delta,
\end{equation}
%
e, portanto,
%
\begin{equation}
    \boxed{
    \gamma=2\Delta+\alpha-2.
    }
    \label{eq:gamma}
\end{equation}
%
\subsection*{(iii) Expoente crítico associado ao calor específico a volume constante}
% 
A energia livre de Gibbs \(g(T,p)\) tem como variáveis naturais \((T,p)\).
Por isso, derivadas segundas de \(g\) em relação a \(T\) a \(p\) fixo dão
diretamente o calor específico a pressão constante, \(c_p\). Entretanto, o
calor específico a volume constante,
%
\begin{equation}
    c_v
    =
    T\left(\frac{\partial s}{\partial T}\right)_v,
    \label{eq:cv_def_p1}
\end{equation}
%
envolve uma derivada mantendo \(v\) fixo, e \(v\) \emph{não} é uma variável
natural de \(g(T,p)\). Não podemos, portanto, obter \(c_v\) simplesmente
tomando \(\partial^2 g/\partial T^2\) a \(v\) constante sem antes mudar de
variáveis. Para contornar isso, usamos uma identidade termodinâmica que
relaciona \(c_p\), \(c_v\), \(\alpha_T\) e \(\kappa_T\), grandezas que podem
ser calculadas diretamente a partir de \(g(T,p)\). A diferença entre os calores específicos é dada por
%
\begin{equation}
    c_p-c_v
    =
    \frac{T\,v\,\alpha_T^{2}}{\kappa_T},
    \label{eq:cp_cv}
\end{equation}
%
onde
%
\begin{equation}
    \alpha_T
    =
    \frac{1}{v}\left(\frac{\partial v}{\partial T}\right)_p,
    \qquad
    \kappa_T
    =
    -\frac{1}{v}\left(\frac{\partial v}{\partial p}\right)_T.
\end{equation}
%
Essa identidade é exata e vale para qualquer sistema termodinâmico simples.
A vantagem é que \(c_p\), \(\alpha_T\) e \(\kappa_T\) são todos obtidos de
derivadas de \(g(T,p)\), de modo que podemos inserir diretamente seus
comportamentos críticos. Do item (iv), o calor específico a pressão constante escala como
%
\begin{equation}
    c_p\sim t^{-\alpha}.
    \label{eq:cp_iii_p1}
\end{equation}
%
Do item (v), o coeficiente de dilatação térmica satisfaz
%
\begin{equation}
    \alpha_T\sim t^{\,1-\alpha-\Delta}
    \quad\Longrightarrow\quad
    \alpha_T^{2}\sim t^{\,2-2\alpha-2\Delta}.
    \label{eq:alphaT2_iii_p1}
\end{equation}
%
Do item (ii), a compressibilidade isotérmica satisfaz
%
\begin{equation}
    \kappa_T\sim t^{\,2-\alpha-2\Delta}.
    \label{eq:kappa_iii_p1}
\end{equation}
%
Dividindo \eqref{eq:alphaT2_iii_p1} por \eqref{eq:kappa_iii_p1},
%
\begin{align}
    \frac{\alpha_T^{2}}{\kappa_T}
    &\sim
    \frac{t^{\,2-2\alpha-2\Delta}}{t^{\,2-\alpha-2\Delta}} \nonumber\\
    &=
    t^{\,(2-2\alpha-2\Delta)-(2-\alpha-2\Delta)} \nonumber\\
    &=
    t^{\,2-2\alpha-2\Delta-2+\alpha+2\Delta} \nonumber\\
    &=
    t^{\,-\alpha}.
    \label{eq:razao_iii_p1}
\end{align}
%
Substituindo \eqref{eq:razao_iii_p1} em \eqref{eq:cp_cv}, e observando que
\(T\to T_c\) e \(v\to v_c\) são finitos e não nulos próximo ao ponto crítico,
%
\begin{equation}
    c_p-c_v
    =
    \frac{T\,v\,\alpha_T^{2}}{\kappa_T}
    \sim
    T_c\,v_c\,
    t^{-\alpha}
    \;\propto\;
    t^{-\alpha}.
    \label{eq:cpminuscv_scaling}
\end{equation}
%
Ou seja, a diferença entre os dois calores específicos tem a mesma potência
crítica que \(c_p\).
%
\begin{equation}
    c_v
    =
    c_p
    -
    \frac{T\,v\,\alpha_T^{2}}{\kappa_T}.
    \label{eq:cv_from_cp_p1}
\end{equation}
%
Ambos os termos do lado direito escalam como \(t^{-\alpha}\): o primeiro por
\eqref{eq:cp_iii_p1}, o segundo por \eqref{eq:cpminuscv_scaling}. Assim, a
contribuição singular dominante de \(c_v\) também escala como \(t^{-\alpha}\):
%
\begin{equation}
    \boxed{
    c_v
    \sim
    t^{\,-\alpha}.
    }
    \label{eq:cv_scaling}
\end{equation}
%
Comparando com a forma canônica \(c_v\sim t^{-\alpha_{c_v}}\), identificamos
%
\begin{equation}
    \boxed{
    \alpha_{c_v}=\alpha.
    }
    \label{eq:expoente_cv_p1}
\end{equation}
%
Esse resultado é
consistente com a relação de Rushbrooke
%
\begin{equation}
    \alpha+2\beta+\gamma=2,
\end{equation}
%
\subsection*{(iv) Expoente crítico associado ao calor específico a pressão constante}
%
O calor específico a pressão constante é definido por
%
\begin{equation}
    c_p
    =T\left(\frac{\partial s}{\partial T}\right)_p
    =-T\left(\frac{\partial^2 g}{\partial T^2}\right)_p.
    \label{eq:cp_def_p1}
\end{equation}
%
Como para uma função \(f(t)\) qualquer vale a regra da cadeia
%
\begin{equation}
    \frac{\partial f}{\partial T}
    =
    f'(t)\,\frac{\partial t}{\partial T}
    =
    \frac{f'(t)}{T_c},
    \label{eq:df_dT_p1}
\end{equation}
%
e, derivando mais uma vez,
%
\begin{equation}
    \frac{\partial^2 f}{\partial T^2}
    =
    \frac{1}{T_c}\,
    \frac{\partial f'(t)}{\partial T}
    =
    \frac{1}{T_c}\,
    f''(t)\,\frac{\partial t}{\partial T}
    =
    \frac{f''(t)}{T_c^{2}}.
    \label{eq:d2f_dT2_p1}
\end{equation}
%
Em particular, para \(f(t)=t^{\,2-\alpha}\), temos
%
\begin{align}
    f'(t)
    &=
    (2-\alpha)\,t^{\,1-\alpha},
    \label{eq:f1_p1}\\
    f''(t)
    &=
    (2-\alpha)(1-\alpha)\,t^{\,-\alpha}.
    \label{eq:f2_p1}
\end{align}
%
Então, para a segunda derivada da energia livre singular ao longo de \(p=p_c\), a parte singular da energia livre de Gibbs é
%
\begin{equation}
    g_s(t,0)\sim F_0\,t^{\,2-\alpha},
    \label{eq:gs_pc_p1}
\end{equation}
%
com \(F_0\equiv F(0)\). Aplicando \eqref{eq:d2f_dT2_p1} com
\(f(t)=F_0\,t^{\,2-\alpha}\),
%
\begin{align}
    \frac{\partial^2 g_s}{\partial T^2}
    &\sim
    \frac{1}{T_c^{2}}\,
    \frac{d^{2}}{dt^{2}}
    \left(F_0\,t^{\,2-\alpha}\right) \nonumber\\
    &=
    \frac{F_0}{T_c^{2}}\,
    (2-\alpha)(1-\alpha)\,
    t^{\,-\alpha}.
    \label{eq:d2gs_dT2_p1}
\end{align}
%
Então para o calor específico a pressão constante, substituindo \eqref{eq:d2gs_dT2_p1} em \eqref{eq:cp_def_p1},
%
\begin{align}
    c_{p,s}
    &\sim
    -\,T\,
    \frac{F_0}{T_c^{2}}\,
    (2-\alpha)(1-\alpha)\,
    t^{\,-\alpha} \nonumber\\
    &=
    -\,
    \frac{T}{T_c^{2}}\,
    (2-\alpha)(1-\alpha)\,F_0\,
    t^{\,-\alpha}.
    \label{eq:cps_prefator_p1}
\end{align}
%
No limite crítico, \(T\to T_c\), de modo que o prefator
%
\begin{equation}
    -\,
    \frac{T}{T_c^{2}}\,
    (2-\alpha)(1-\alpha)\,F_0
    \;\xrightarrow[T\to T_c]{}\;
    -\,
    \frac{(2-\alpha)(1-\alpha)\,F_0}{T_c},
    \label{eq:prefator_finito_p1}
\end{equation}
%
é uma constante finita e não nula (desde que \(F_0\neq0\) e
\(\alpha\neq1,2\)). Logo, a parte singular do calor específico a pressão
constante escala como
%
\begin{equation}
    \boxed{
    c_p
    \sim
    t^{\,-\alpha}.
    }
    \label{eq:cp_scaling}
\end{equation}
%
Comparando com a forma canônica \(c_p\sim t^{-\alpha_{c_p}}\), identificamos
diretamente
%
\begin{equation}
    \boxed{
    \alpha_{c_p}=\alpha.
    }
    \label{eq:expoente_cp_p1}
\end{equation}
%
Ou seja, o expoente crítico associado ao calor específico a pressão
constante é justamente o expoente \(\alpha\) que aparece na hipótese de
escala da energia livre \eqref{eq:g_scaling}.
%
\subsection*{(v) Expoente crítico associado ao coeficiente de dilatação térmica}
% 
Definimos o coeficiente de dilatação térmica por
%
\begin{equation}
    \alpha_T
    =\frac{1}{v}\left(\frac{\partial v}{\partial T}\right)_p
    =\frac{1}{v}
    \frac{\partial^2g}{\partial T\,\partial p}.
\end{equation}
Usando \eqref{eq:d2gpt}, obtemos no caminho \(p=p_c\)
\begin{equation}
    \alpha_{T,s}
    \sim
    \frac{2-\alpha-\Delta}{v_cp_cT_c}
    F_1t^{1-\alpha-\Delta}.
\end{equation}
Portanto,
\begin{equation}
    \boxed{
    \alpha_T\sim t^{1-\alpha-\Delta}.
    }
    \label{eq:alphaT_scaling}
\end{equation}
Se escrevermos como lei de potência divergente, \(\alpha_T\sim t^{-\phi}\),
então
\begin{equation}
    \boxed{
    \phi=\alpha+\Delta-1.
    }
    \label{eq:phi}
\end{equation}
Além disso, como \(v_s\sim t^{2-\alpha-\Delta}\), identificando a parte
singular de \(v\) com o expoente de ordem \(\beta\),
\begin{equation}
    \boxed{
    \beta=2-\alpha-\Delta.
    }
    \label{eq:beta}
\end{equation}
Assim, o resultado para a dilatação pode ser escrito também como
\begin{equation}
    \boxed{
    \alpha_T\sim t^{\beta-1}=t^{-(1-\beta)}.
    }
\end{equation}
% ------------------------------------------------------------
%
\begin{center}
\renewcommand{\arraystretch}{1.5}
\begin{tabular}{>{\centering\arraybackslash}m{0.30\linewidth}
                >{\centering\arraybackslash}m{0.28\linewidth}
                >{\centering\arraybackslash}m{0.28\linewidth}}
\toprule
\textbf{Grandeza} & \textbf{Comportamento crítico} & \textbf{Expoente}\\
\midrule
Entropia 
    & \(s_s\sim t^{1-\alpha}\)
    & \(1-\alpha\) \\
Compressibilidade isotérmica
    & \(\kappa_T\sim t^{-\gamma}\)
    & \(\boxed{\gamma=2\Delta+\alpha-2}\) \\
Calor específico a \(v\) constante
    & \(c_v\sim t^{-\alpha}\)
    & \(\boxed{\alpha}\) \\
Calor específico a \(p\) constante
    & \(c_p\sim t^{-\alpha}\)
    & \(\boxed{\alpha}\) \\
Dilatação térmica
    & \(\alpha_T\sim t^{1-\alpha-\Delta}\)
    & \(\boxed{\phi=\alpha+\Delta-1}\) \\
\bottomrule
\end{tabular}
\end{center}
% ============================================================
\section*{Problema 2 -- Verificação da forma de escala para o níquel}
% ============================================================

Os dados de Weiss e Forrer para a magnetização do níquel, analisados por
Kouvel e colaboradores e corrigidos para efeitos de desmagnetização,
podem ser testados à luz da hipótese de escala próxima do ponto crítico.
O enunciado fornece os parâmetros críticos
%
\begin{equation}
    T_c = 627.4\,\mathrm{K},
    \qquad
    \beta = 0.378(4),
    \qquad
    \gamma = 1.34(1),
    \qquad
    \delta = 4.58(5),
\end{equation}
%
e propõe a forma de escala
%
\begin{equation}
    m
    =
    |t|^\beta
    Y_{\pm}
    \left(
        \frac{H}{|t|^\Delta}
    \right),
    \label{eq:escala_Ni}
\end{equation}
%
onde
%
\begin{equation}
    t
    =
    \frac{T-T_c}{T_c},
    \qquad
    \Delta=\beta+\gamma,
    \label{eq:variaveis_Ni_1}
\end{equation}
%
e $Y_-$ e $Y_+$ representam, respectivamente, as funções de escala
abaixo e acima da temperatura crítica. O campo magnético que deve ser utilizado na análise é o campo corrigido pelos efeitos de
desmagnetização,
%
\begin{equation}
    H = H_0 - 39.3\,m.
    \label{eq:H_corrigido_Ni}
\end{equation}
%
A partir dos valores experimentais de $\beta$ e $\gamma$, obtemos
%
\begin{equation}
    \Delta
    =
    \beta+\gamma
    =
    0.378+1.34
    =
    1.718.
    \label{eq:Delta_Ni}
\end{equation}
%
Se desejarmos propagar apenas as incertezas estatísticas fornecidas,
assumindo-as independentes, temos
%
\begin{equation}
    \sigma_\Delta
    =
    \sqrt{\sigma_\beta^2+\sigma_\gamma^2}
    =
    \sqrt{(0.004)^2+(0.01)^2}
    \simeq0.011,
\end{equation}
%
de modo que
%
\begin{equation}
    \Delta=1.718(11).
\end{equation}
%
Também podemos verificar a relação de escala de Widom,
%
\begin{equation}
    \gamma=\beta(\delta-1).
    \label{eq:Widom_Ni}
\end{equation}
%
Usando os valores experimentais,
%
\begin{equation}
    \beta(\delta-1)
    =
    0.378(4)\,[4.58(5)-1]
    =
    1.353.
\end{equation}
%
A incerteza associada a esse valor é aproximadamente
%
\begin{equation}
    \sigma_{\beta(\delta-1)}
    =
    \sqrt{
        (\delta-1)^2\sigma_\beta^2
        +
        \beta^2\sigma_\delta^2
    }
    \simeq0.017,
\end{equation}
%
resultando em
%
\begin{equation}
    \beta(\delta-1)=1.353(17).
\end{equation}
%
Esse resultado é compatível com o valor experimental $\gamma=1.34(1)$, dentro das incertezas. Definimos as duas variáveis adimensionais
%
\begin{equation}
    x
    \equiv
    \frac{H}{|t|^\Delta},
    \qquad
    y
    \equiv
    \frac{m}{|t|^\beta}.
    \label{eq:xy_Ni}
\end{equation}
%
A Eq.~\eqref{eq:escala_Ni} pode então ser re-escrita simplesmente como
%
\begin{equation}
    y=Y_{\pm}(x).
    \label{eq:collapse_function_Ni}
\end{equation}
%
Portanto, se a hipótese de escala for válida, medidas realizadas em
diferentes temperaturas suficientemente próximas de $T_c$ devem,
quando expressas em termos de $(x,y)$, colapsar em uma mesma função de
escala para cada lado da transição. Para $T<T_c$, teremos a função $Y_-(x)$, enquanto para $T>T_c$ teremos
a função $Y_+(x)$.
%
As temperaturas da tabela original são fornecidas em graus Celsius.
Devemos convertê-las para Kelvin segundo
%
\begin{equation}
    T(\mathrm{K})
    =
    T(^\circ\mathrm{C})+273.15.
\end{equation}
%
Assim,
%
\begin{align}
    348.78^\circ\mathrm{C}
    &\longrightarrow
    621.93\,\mathrm{K},
    \\
    350.66^\circ\mathrm{C}
    &\longrightarrow
    623.81\,\mathrm{K},
    \\
    352.53^\circ\mathrm{C}
    &\longrightarrow
    625.68\,\mathrm{K},
    \\
    358.18^\circ\mathrm{C}
    &\longrightarrow
    631.33\,\mathrm{K}.
\end{align}
%
Usando $T_c=627.4\,\mathrm K$, encontramos
%
\begin{equation}
    t=\frac{T-T_c}{T_c},
\end{equation}
%
e, portanto,
%
\begin{align}
    T=348.78^\circ\mathrm C:
    &&t&=-0.008719,
    \\
    T=350.66^\circ\mathrm C:
    &&t&=-0.005722,
    \\
    T=352.53^\circ\mathrm C:
    &&t&=-0.002741,
    \\
    T=358.18^\circ\mathrm C:
    &&t&=+0.006264.
\end{align}
%
As três primeiras temperaturas estão abaixo de $T_c$, enquanto
$358.18^\circ\mathrm C$ está acima de $T_c$.
%
Para cada ponto experimental, o campo que deve ser utilizado na
hipótese de escala não é diretamente $H_0$, mas o campo corrigido dado pela equação~\eqref{eq:H_corrigido_Ni}.
%
Por exemplo, para
$T=348.78^\circ\mathrm C$,
$H_0=1355\,\mathrm G$ e $m=14.307$, temos
%
\begin{align}
    H
    &=
    1355-39.3(14.307)
    \nonumber\\
    &=
    792.735\,\mathrm G.
\end{align}
%
Como
%
\begin{equation}
    |t|^\Delta
    =
    |{-0.008719}|^{1.718}
    \simeq2.895\times10^{-4},
\end{equation}
%
segue que
%
\begin{equation}
    x
    =
    \frac{792.735}{2.895\times10^{-4}}
    \simeq2.738\times10^6.
\end{equation}
%
Analogamente,
%
\begin{equation}
    |t|^\beta
    =
    |{-0.008719}|^{0.378}
    \simeq0.1665,
\end{equation}
%
e portanto
%
\begin{equation}
    y
    =
    \frac{14.307}{0.1665}
    \simeq85.91.
\end{equation}
%
Esse procedimento é aplicado a todos os pontos da tabela. A tabela abaixo apresenta os dados experimentais corrigidos e as correspondentes variáveis de escala.
%
\begin{table}[H]
\centering
\small
\caption{
Dados de magnetização do níquel corrigidos pelos efeitos de
desmagnetização e variáveis de escala. Foram utilizados
$T_c=627.4\,\mathrm K$, $\beta=0.378$ e $\Delta=1.718$.
}
\label{tab:dados_Ni}
\begin{tabular}{@{}ccccc@{}}
\toprule
$H_0\,(\mathrm G)$
&
$m$
&
$H\,(\mathrm G)$
&
$x=H/|t|^\Delta$
&
$y=m/|t|^\beta$
\\
\midrule

\multicolumn{5}{c}{
$T=348.78^\circ\mathrm C$,
$t=-0.008719$,
$|t|^\beta=0.1665$,
$|t|^\Delta=2.895\times10^{-4}$
}
\\

430
& 11.361
& -16.487
& $-5.695\times10^4$
& 68.223
\\

890
& 13.787
& 348.171
& $1.203\times10^6$
& 82.791
\\

1355
& 14.307
& 792.735
& $2.738\times10^6$
& 85.913
\\

3230
& 15.651
& 2614.916
& $9.032\times10^6$
& 93.984
\\

6015
& 16.938
& 5349.337
& $1.848\times10^7$
& 101.713
\\

10070
& 18.209
& 9354.386
& $3.231\times10^7$
& 109.345
\\

14210
& 19.212
& 13454.968
& $4.647\times10^7$
& 115.368
\\

17775
& 19.976
& 16989.943
& $5.868\times10^7$
& 119.956
\\

\midrule

\multicolumn{5}{c}{
$T=350.66^\circ\mathrm C$,
$t=-0.005722$,
$|t|^\beta=0.1420$,
$|t|^\Delta=1.404\times10^{-4}$
}
\\

430
& 10.452
& 19.236
& $1.370\times10^5$
& 73.595
\\

1355
& 12.646
& 858.012
& $6.110\times10^6$
& 89.043
\\

3230
& 14.322
& 2667.145
& $1.899\times10^7$
& 100.844
\\

6015
& 15.780
& 5394.846
& $3.842\times10^7$
& 111.110
\\

10070
& 17.271
& 9391.250
& $6.687\times10^7$
& 121.608
\\

14210
& 18.348
& 13488.924
& $9.605\times10^7$
& 129.192
\\

17775
& 19.141
& 17022.759
& $1.212\times10^8$
& 134.775
\\

\midrule

\multicolumn{5}{c}{
$T=352.53^\circ\mathrm C$,
$t=-0.002741$,
$|t|^\beta=0.1075$,
$|t|^\Delta=3.967\times10^{-5}$
}
\\

430
& 8.114
& 111.120
& $2.801\times10^6$
& 75.453
\\

1355
& 10.626
& 937.398
& $2.363\times10^7$
& 98.812
\\

3230
& 12.752
& 2728.846
& $6.879\times10^7$
& 118.582
\\

6015
& 14.519
& 5444.403
& $1.372\times10^8$
& 135.014
\\

10070
& 16.198
& 9433.419
& $2.378\times10^8$
& 150.627
\\

14210
& 17.395
& 13526.377
& $3.410\times10^8$
& 161.758
\\

17775
& 18.293
& 17056.085
& $4.300\times10^8$
& 170.109
\\

\midrule

\multicolumn{5}{c}{
$T=358.18^\circ\mathrm C$,
$t=+0.006264$,
$|t|^\beta=0.1470$,
$|t|^\Delta=1.641\times10^{-4}$
}
\\

1355
& 4.235
& 1188.564
& $7.245\times10^6$
& 28.817
\\

3230
& 7.567
& 2932.617
& $1.788\times10^7$
& 51.489
\\

6015
& 10.312
& 5609.738
& $3.419\times10^7$
& 70.167
\\

10070
& 12.775
& 9567.942
& $5.832\times10^7$
& 86.927
\\

14210
& 14.385
& 13644.669
& $8.317\times10^7$
& 97.882
\\

17775
& 15.504
& 17165.693
& $1.046\times10^8$
& 105.496
\\

\bottomrule
\end{tabular}
\end{table}
%
Para testar a hipótese de escala, construímos o gráfico de
%
\begin{equation}
    y=\frac{m}{|t|^\beta}
\end{equation}
%
em função de
%
\begin{equation}
    x=\frac{H}{|t|^\Delta}.
\end{equation}
%
Para os três conjuntos de dados abaixo de $T_c$, esperamos que os
pontos se aproximem de uma única função $Y_-(x)$. Para o conjunto acima
de $T_c$, esperamos a função $Y_+(x)$. Há, entretanto, um detalhe importante. Para
$T=348.78^\circ\mathrm C$ e $H_0=430\,\mathrm G$, a correção de
desmagnetização fornece
%
\begin{equation}
    H
    =
    430-39.3(11.361)
    =
    -16.487\,\mathrm G,
\end{equation}
%
e, portanto,
%
\begin{equation}
    x<0.
\end{equation}
%
Esse ponto não pode ser representado em um gráfico log-log, uma vez que
o logaritmo de $x$ não está definido para $x<0$. Portanto, esse único
ponto será omitido da figura de colapso em escala log-log.
%
\begin{figure}[H]
\centering
\begin{tikzpicture}
\begin{loglogaxis}[
    width=0.78\linewidth,
    height=0.58\linewidth,
    xlabel={\(x=H/|t|^{\Delta}\)},
    ylabel={\(y=m/|t|^{\beta}\)},
    grid=both,
    legend style={
        at={(0.02,0.98)},
        anchor=north west,
        font=\small
    },
    legend cell align=left,
    mark size=1.7pt,
    xmin=1e5,
    xmax=5e8,
    ymin=20,
    ymax=200,
]

% ------------------------------------------------------------
% T < Tc
% ------------------------------------------------------------
\addplot+[
    only marks,
    mark=*,
    blue
] coordinates {
    (1.203e6,82.791)
    (2.738e6,85.913)
    (9.032e6,93.984)
    (1.848e7,101.713)
    (3.231e7,109.345)
    (4.647e7,115.368)
    (5.868e7,119.956)

    (1.370e5,73.595)
    (6.110e6,89.043)
    (1.899e7,100.844)
    (3.842e7,111.110)
    (6.687e7,121.608)
    (9.605e7,129.192)
    (1.212e8,134.775)

    (2.801e6,75.453)
    (2.363e7,98.812)
    (6.879e7,118.582)
    (1.372e8,135.014)
    (2.378e8,150.627)
    (3.410e8,161.758)
    (4.300e8,170.109)
};
\addlegendentry{\(T<T_c\) (\(Y_-\))}

% ------------------------------------------------------------
% T > Tc
% ------------------------------------------------------------
\addplot+[
    only marks,
    mark=square*,
    red
] coordinates {
    (7.245e6,28.817)
    (1.788e7,51.489)
    (3.419e7,70.167)
    (5.832e7,86.927)
    (8.317e7,97.882)
    (1.046e8,105.496)
};
\addlegendentry{\(T>T_c\) (\(Y_+\))}
\end{loglogaxis}
\end{tikzpicture}
%
\caption{
Teste da forma de escala da magnetização do níquel.
Os dados das três temperaturas abaixo de $T_c$ (pontos azuis) apresentam um colapso
aproximado sobre uma única ramificação $Y_-(x)$ quando expressos nas
variáveis reduzidas
$x=H/|t|^\Delta$ e $y=m/|t|^\beta$.
Os pontos disponíveis acima de $T_c$ (quadrados vermelhos) formam uma segunda ramificação
$Y_+(x)$. O ponto correspondente a
$T=348.78^\circ\mathrm C$ e $H_0=430\,\mathrm G$ possui $x<0$ e, por
isso, não é mostrado na representação log-log.
}
\label{fig:colapso_Ni}
\end{figure}
%
A figura mostra que os três conjuntos de dados obtidos para
$T<T_c$ podem ser aproximadamente reunidos em uma mesma curva quando
as variáveis originais são reescaladas. Podemos observar, por exemplo, que em regiões de valores semelhantes de $x$, os valores de $y$ obtidos em temperaturas diferentes tendem a se aproximar. Isso constitui evidência de que a dependência explícita em $t$ pode ser removida pela combinação de variáveis prevista pela teoria de escala. Os pontos obtidos para $T=358.18^\circ\mathrm C>T_c$ formam uma ramificação distinta, compatível com a existência de uma função
$Y_+(x)$ diferente de $Y_-(x)$. Entretanto, é importante observar que há somente \emph{uma} temperatura acima de $T_c$ no conjunto de dados fornecido. Assim, esses pontos demonstram a forma de uma segunda ramificação, mas não permitem um teste independente de colapso de $Y_+$, pois não há uma segunda temperatura acima de $T_c$ com a qual realizar a comparação. Por outro lado, as três temperaturas abaixo de $T_c$ permitem um teste mais direto da hipótese de escala, pois diferentes valores de $t$
devem produzir a mesma função $Y_-(x)$ após o reescalonamento. Os dados experimentais do níquel são compatíveis com a forma de escala
% ============================================================
\section*{Problema 3 -- Magnetização espontânea do modelo de Ising na rede quadrada}
% ============================================================
A magnetização espontânea do modelo de Ising ferromagnético com interações de
primeiros vizinhos na rede quadrada é dada pela fórmula exata de Onsager,
%
\begin{equation}
    m = \left\{1 - \left[\sinh(2K)\right]^{-4}\right\}^{1/8},
    \qquad
    K = \beta J = \frac{J}{k_B T}.
    \label{eq:onsager_m}
\end{equation}
%
Desejamos verificar a forma assintótica
%
\begin{equation}
    m \sim B\,(-t)^{\beta},
    \qquad
    t = \frac{T-T_c}{T_c},
    \label{eq:assint_m}
\end{equation}
%
nas vizinhanças do ponto crítico, e determinar o prefator \(B\) e o expoente crítico \(\beta\). A transição de fase ocorre quando o argumento da raiz se anula, isto é,
%
\begin{equation}
    \sinh(2K_c) = 1.
    \label{eq:critical_condition}
\end{equation}
%
A função $\operatorname{arcsinh}(x)$ pode ser escrita usando logaritmos naturais com a fórmula:
%
\[
\operatorname{arcsinh}(x) = \ln \left(x + \sqrt{x^{2} + 1}\right)
\]
%
Substituindo $x = 1$:
%
\[
\operatorname{arcsinh}(1) = \ln \left(1 + \sqrt{1^{2} + 1}\right) = \ln (1 + \sqrt{2})
\]
%
de modo que, resolvendo para \(K_c\),
%
\begin{equation}
    2K_c = \operatorname{arcsinh}(1) = \ln(1+\sqrt{2}),
\end{equation}
%
\begin{equation}
    K_c = \frac{J}{k_B T_c}
    = \frac{1}{2}\ln(1+\sqrt{2})
    \approx 0.4406868.
    \label{eq:Kc}
\end{equation}
% 
Escrevemos a temperatura como
%
\begin{equation}
    T = T_c(1+t),
    \qquad
    t = \frac{T-T_c}{T_c}.
\end{equation}
%
Como \(K = J/(k_B T)\), temos
%
\begin{equation}
    K = \frac{J}{k_B T_c(1+t)}
    = \frac{K_c}{1+t}
    \approx K_c\left(1 - t + \mathcal O(t^2)\right).
    \label{eq:K_expansion}
\end{equation}
%
Portanto,
%
\begin{equation}
    K - K_c \approx -K_c\,t.
    \label{eq:K_difference}
\end{equation}
%
Definimos
%
\begin{equation}
    f(K) \equiv \sinh(2K).
\end{equation}
%
No ponto crítico,
%
\begin{equation}
    f(K_c) = \sinh(2K_c) = 1.
\end{equation}
%
A derivada é
%
\begin{equation}
    f'(K) = 2\cosh(2K).
\end{equation}
%
Usando \(\cosh^2(2K_c) = 1 + \sinh^2(2K_c) = 2\), obtemos
%
\begin{equation}
    f'(K_c) = 2\sqrt{2}.
\end{equation}
%
Assim, a expansão de Taylor de \(f(K)\) em torno de \(K_c\) fornece
%
\begin{align}
    \sinh(2K)
    &= f(K_c) + f'(K_c)(K-K_c) + \mathcal O\big((K-K_c)^2\big) \nonumber\\
    &\approx 1 + 2\sqrt{2}\,(K-K_c) \nonumber\\
    &\approx 1 - 2\sqrt{2}\,K_c\,t.
    \label{eq:sinh_expansion}
\end{align}
%
Definimos a constante auxiliar
%
\begin{equation}
    a \equiv 2\sqrt{2}\,K_c,
    \label{eq:a_def}
\end{equation}
%
de modo que
%
\begin{equation}
    \sinh(2K) \approx 1 - a\,t.
    \label{eq:sinh_linear}
\end{equation}
%
Substituindo \eqref{eq:sinh_linear} na fórmula de Onsager \eqref{eq:onsager_m},
\begin{equation}
    \left[\sinh(2K)\right]^{-4}
    \approx (1 - a\,t)^{-4}
    \approx 1 + 4a\,t + \mathcal O(t^2),
    \label{eq:sinh_power}
\end{equation}
%
onde usamos a expansão binomial \((1-x)^{-4} \approx 1 + 4x\) para \(x \ll 1\). Portanto,
%
\begin{equation}
    1 - \left[\sinh(2K)\right]^{-4}
    \approx -4a\,t
    = 4a\,(-t).
    \label{eq:argument_m}
\end{equation}
%
Como a magnetização espontânea é não nula apenas para \(T < T_c\), isto é,
\(t < 0\), temos \(-t = |t| > 0\). Assim,
%
\begin{equation}
    m = \left\{1 - \left[\sinh(2K)\right]^{-4}\right\}^{1/8}
    \approx \left[4a\,(-t)\right]^{1/8}
    = (4a)^{1/8}\,(-t)^{1/8}.
    \label{eq:m_scaling}
\end{equation}
%
Comparando com a forma assintótica \eqref{eq:assint_m}, identificamos
%
\begin{equation}
    \boxed{
    \beta = \frac{1}{8}
    }
    \label{eq:beta_onsager}
\end{equation}
%
O prefator \(B\) é
%
\begin{equation}
    B = (4a)^{1/8} = \left(8\sqrt{2}\,K_c\right)^{1/8}.
    \label{eq:B_Kc}
\end{equation}
%
Usando o valor explícito de \(K_c\) dado em \eqref{eq:Kc},
%
\begin{equation}
    \boxed{
    B = \left[8\sqrt{2}\cdot \frac{1}{2}\ln(1+\sqrt{2})\right]^{1/8}
    = \left[4\sqrt{2}\,\ln(1+\sqrt{2})\right]^{1/8}
    \approx 1.222.
    }
    \label{eq:B_final}
\end{equation}
%
Nas vizinhanças do ponto crítico, a magnetização espontânea do modelo de Ising na rede quadrada obedece à forma assintótica
%
\begin{equation}
    \boxed{
    m \sim B\,(-t)^{\beta},
    \qquad
    \beta = \frac{1}{8},
    \qquad
    B = \left[4\sqrt{2}\,\ln(1+\sqrt{2})\right]^{1/8} \approx 1.222,
    }
    \label{eq:resultado_final}
\end{equation}
%
para \(t < 0\). Para \(t > 0\), a magnetização espontânea é nula, \(m = 0\). O expoente \(\beta = 1/8\) é o valor exato da classe de universalidade do modelo de Ising bidimensional, diferindo substancialmente do valor clássico de campo médio \(\beta = 1/2\) que calculamos na primeira série de exercícios.
% ============================================================
\section*{Problema 4 -- Expansão de altas temperaturas do modelo de Ising na rede quadrada}
% 
Considere o modelo de Ising ferromagnético na rede quadrada simples, com interações de primeiros vizinhos e campo externo nulo. Adotamos o hamiltoniano
com o sinal que torna a interação ferromagnética,
%
\begin{equation}
    \mathcal H = -J \sum_{\langle i,j\rangle} \sigma_i \sigma_j,
    \qquad J>0,\quad \sigma_i=\pm1,
    \label{eq:H_ising_square}
\end{equation}
%
onde a soma percorre todos os pares de primeiros vizinhos. A função canônica de partição é
%
\begin{equation}
    Z = \sum_{\{\sigma_i\}} \e{-\beta \mathcal H}
    = \sum_{\{\sigma_i\}} \exp\left(K \sum_{\langle i,j\rangle} \sigma_i \sigma_j\right),
    \qquad K=\beta J.
    \label{eq:Z_def}
\end{equation}
%
Na rede quadrada com $N$ sítios e condições periódicas de contorno, o número de
ligações (pares de primeiros vizinhos) é $2N$.
% ------------------------------------------------------------
\subsection*{(a) Expansão de altas temperaturas e grafos associados}
% ------------------------------------------------------------
A variável $x \equiv \sigma_i \sigma_j$ assume apenas dois valores discretos possíveis (para um sistema de spin-$1/2$), $x \in \{-1, +1\}$, dado que $\sigma_i = \pm 1$ e $\sigma_j = \pm 1$. Por conseguinte, para qualquer número inteiro $n \ge 0$, vale a propriedade binária:
%
\begin{equation}
    x^n = 
    \begin{cases} 
    1, & \text{se } n \text{ é par}, \\
    x, & \text{se } n \text{ é ímpar}.
    \end{cases}
\end{equation}
%
Expandindo a função exponencial $\e{K x}$ em série de Taylor em torno de $x=0$:
%
\begin{equation}
    \e{K x} = \sum_{n=0}^{\infty} \frac{(K x)^n}{n!} = \sum_{n \text{ par}} \frac{K^n x^n}{n!} + \sum_{n \text{ ímpar}} \frac{K^n x^n}{n!}.
\end{equation}
%
Aplicando a propriedade das potências de $x$, escrevemos $x^{2m} = 1$ para termos pares e $x^{2m+1} = x$ para termos ímpares, colocando $x$ em evidência no segundo somatório:
%
\begin{equation}
    \e{K x} = \left( \sum_{m=0}^{\infty} \frac{K^{2m}}{(2m)!} \right) + x \left( \sum_{m=0}^{\infty} \frac{K^{2m+1}}{(2m+1)!} \right).
\end{equation}
%
Reconhecendo as definições em série de Taylor do cosseno hiperbólico e do seno hiperbólico:
%
\begin{equation}
    \cosh K = \sum_{m=0}^{\infty} \frac{K^{2m}}{(2m)!}, \qquad \sinh K = \sum_{m=0}^{\infty} \frac{K^{2m+1}}{(2m+1)!},
\end{equation}
%
obtemos a identidade linear em $x$:
%
\begin{equation}
    \e{K \sigma_i \sigma_j} = \cosh K + \sigma_i \sigma_j \sinh K.
\end{equation}
%
Fatorando $\cosh K$ da expressão obtida:
%
\begin{equation}
    \e{K \sigma_i \sigma_j} = \cosh K \left( 1 + \frac{\sinh K}{\cosh K} \sigma_i \sigma_j \right) = \cosh K \left(1 + t\, \sigma_i \sigma_j\right),
    \label{eq:identity}
\end{equation}
%
onde definiu-se a variável de expansão de altas temperaturas $t \equiv \tanh K$. Agora, como a função de partição canônica do sistema é definida por \eqref{eq:Z_def}, usamos a propriedade fundamental da exponencial que converte somas em produtos:
%
\begin{equation}
    Z = \sum_{\{\sigma_i\}} \prod_{\langle i,j\rangle} \e{K \sigma_i \sigma_j}.
\end{equation}
%
Substituindo a identidade \eqref{eq:identity} derivada anteriormente para cada ligação $\langle i,j \rangle$:
%
\begin{equation}
    Z = \sum_{\{\sigma_i\}} \prod_{\langle i,j\rangle} \left[ \cosh K \left(1 + t\, \sigma_i \sigma_j\right) \right].
\end{equation}
%
Como o produto opera sobre todas as ligações de primeiros vizinhos, o termo escalar $\cosh K$ é multiplicado tantas vezes quanto o número total de ligações da rede. Na rede quadrada simples com $N$ sítios e número de coordenação $z=4$ com condições periódicas de contorno, o número total de pares de primeiros vizinhos $N_{\text{ligações}}$ é:
%
\begin{equation}
    N_{\text{ligações}} = \frac{z N}{2} = \frac{4 N}{2} = 2N.
\end{equation}
%
Logo:
%
\begin{equation}
    \prod_{\langle i,j\rangle} \cosh K = (\cosh K)^{2N}.
\end{equation}
%
Como $(\cosh K)^{2N}$ não depende das configurações das variáveis de spin $\{\sigma_i\}$, ele pode ser fatorado para fora do somatório sobre as configurações de spin:
%
\begin{equation}
    Z = (\cosh K)^{2N} \sum_{\{\sigma_i\}} \prod_{\langle i,j\rangle} \left(1 + t\, \sigma_i \sigma_j\right),
    \label{eq:Z_prod}
\end{equation}
%
Expandindo o produto,
%
\begin{equation}
    \prod_{\langle i,j\rangle} \left(1 + t\, \sigma_i \sigma_j\right)
    = \sum_{A \subseteq \mathcal B} t^{|A|}
    \prod_{\langle i,j\rangle \in A} \sigma_i \sigma_j,
\end{equation}
%
onde $\mathcal B$ é o conjunto de todas as $2N$ ligações e $A$ é um subconjunto de ligações. Portanto,
%
\begin{equation}
    Z = (\cosh K)^{2N} \sum_{A \subseteq \mathcal B} t^{|A|}
    \sum_{\{\sigma_i\}} \prod_{\langle i,j\rangle \in A} \sigma_i \sigma_j.
    \label{eq:Z_sum_A}
\end{equation}
%
A soma sobre os spins é não nula somente se cada vértice da rede aparece um número par de vezes no subconjunto $A$, isto é, se todos os vértices têm grau par. Nesse caso, a soma vale $2^N$. Logo,
%
\begin{equation}
    Z = (\cosh K)^{2N} 2^N
    \sum_{\substack{A \subseteq \mathcal B \\ \text{grau par}}} t^{|A|}.
    \label{eq:Z_even}
\end{equation}
%
Assim, os coeficientes da expansão em $t$ são dados pelo número de subgrafos com todos os vértices de grau par (subgrafos de Euler) na rede quadrada.

\noindent\textbf{Ordem $t^4$:} O único subgrafo conexo com $4$ ligações e todos os vértices de grau par é um quadrado elementar (plaqueta). Há $N$ quadrados elementares na rede quadrada. Portanto, o coeficiente de $t^4$ é $N$. Ver Figura~\ref{fig:ordem_t4}.

\noindent\textbf{Ordem $t^6$:} O único subgrafo conexo com $6$ ligações e graus pares é um ciclo de comprimento $6$, que na rede quadrada corresponde a um retângulo $1\times 2$ (ou $2\times 1$). Há $N$ retângulos horizontais e $N$ verticais, totalizando $2N$. Logo, o coeficiente de $t^6$ é $2N$. Ver Figura~\ref{fig:ordem_t6_h} e \ref{fig:ordem_t6_v}.
%
\begin{figure}[htbp]
    \centering
    % -------------------------------------------------------------
    % SUBFIGURA (a): ORDEM t^4
    % -------------------------------------------------------------
    \begin{subfigure}[b]{0.31\textwidth}
        \centering
        \begin{tikzpicture}[scale=0.75]
            % Rede quadrada de fundo
            \draw[gray!40, thin, dashed] (-0.3,-0.3) grid (2.3,2.3);
            \foreach \x in {0,...,2} {
                \foreach \y in {0,...,2} {
                    \fill[gray!50] (\x,\y) circle (2pt);
                }
            }
            
            % Plaqueta elementar (ordem t^4)
            \draw[blue!80!black, line width=1.6pt, line join=round] (0,0) -- (1,0) -- (1,1) -- (0,1) -- cycle;
            
            % Vértices ativos
            \foreach \x/\y in {0/0, 1/0, 1/1, 0/1} {
                \fill[blue!80!black] (\x,\y) circle (3pt);
            }
            
            \node at (0.5,0.5) [scale=0.8, font=\bfseries, color=blue!80!black] {Plaqueta};
        \end{tikzpicture}
        \caption{Ordem $t^4$: plaqueta $1\times 1$ ($N$ posições).}
        \label{fig:ordem_t4}
    \end{subfigure}
    \hfill
    % -------------------------------------------------------------
    % SUBFIGURA (b): ORDEM t^6 HORIZONTAL
    % -------------------------------------------------------------
    \begin{subfigure}[b]{0.33\textwidth}
        \centering
        \begin{tikzpicture}[scale=0.75]
            % Rede quadrada de fundo
            \draw[gray!40, thin, dashed] (-0.3,-0.3) grid (3.3,2.3);
            \foreach \x in {0,...,3} {
                \foreach \y in {0,...,2} {
                    \fill[gray!50] (\x,\y) circle (2pt);
                }
            }
            
            % Retângulo horizontal 1x2 (ordem t^6)
            \draw[red!80!black, line width=1.6pt, line join=round] (0,0) -- (2,0) -- (2,1) -- (0,1) -- cycle;
            
            % Vértices ativos
            \foreach \x/\y in {0/0, 1/0, 2/0, 2/1, 1/1, 0/1} {
                \fill[red!80!black] (\x,\y) circle (3pt);
            }
            
            \node at (1,0.5) [scale=0.8, font=\bfseries, color=red!80!black] {Horizontal $1\times 2$};
        \end{tikzpicture}
        \caption{Ordem $t^6$: $1\times 2$ ($N$ posições).}
        \label{fig:ordem_t6_h}
    \end{subfigure}
    \hfill
    % -------------------------------------------------------------
    % SUBFIGURA (c): ORDEM t^6 VERTICAL
    % -------------------------------------------------------------
    \begin{subfigure}[b]{0.31\textwidth}
        \centering
        \begin{tikzpicture}[scale=0.75]
            % Rede quadrada de fundo
            \draw[gray!40, thin, dashed] (-0.3,-0.3) grid (2.3,3.3);
            \foreach \x in {0,...,2} {
                \foreach \y in {0,...,3} {
                    \fill[gray!50] (\x,\y) circle (2pt);
                }
            }
            
            % Retângulo vertical 2x1 (ordem t^6)
            \draw[red!80!black, line width=1.6pt, line join=round] (0,0) -- (1,0) -- (1,2) -- (0,2) -- cycle;
            
            % Vértices ativos
            \foreach \x/\y in {0/0, 1/0, 1/1, 1/2, 0/2, 0/1} {
                \fill[red!80!black] (\x,\y) circle (3pt);
            }
            
            \node at (0.5,1) [scale=0.8, font=\bfseries, color=red!80!black, rotate=90] {Vertical $2\times 1$};
        \end{tikzpicture}
        \caption{Ordem $t^6$: $2\times 1$ ($N$ posições).}
        \label{fig:ordem_t6_v}
    \end{subfigure}

    \caption{Subgrafos de Euler permitidos na expansão de alta temperatura: (a) laço de 4 arestas ($N$ combinações) e (b, c) laços de 6 arestas em ambas as orientações ($N + N = 2N$ combinações).}
    \label{fig:subgrafos_euler}
\end{figure}

\medskip
\noindent\textbf{Ordem $t^8$:} Os subgrafos com $8$ ligações e graus pares podem ser:
%
\begin{itemize}
    \item \textbf{Dois quadrados disjuntos} (dois ciclos de comprimento $4$ sem
    vértices em comum). O número de pares não ordenados de quadrados disjuntos é
    \begin{equation}
        \binom{N}{2} - 4N
        = \frac{N(N-1)}{2} - 4N
        = \frac{N^2}{2} - \frac{9N}{2},
    \end{equation}
    onde subtraímos os $4N$ pares de quadrados que compartilham pelo menos um
    vértice.
    \item \textbf{Ciclos de comprimento $8$}: há $7$ tipos distintos de ciclos de
    comprimento $8$ na rede quadrada, cada um ocorrendo em $N$ posições,
    totalizando $7N$.
    \item \textbf{Dois quadrados que compartilham um único vértice}: em cada
    vértice, há $2$ pares de quadrados que se tocam apenas naquele vértice,
    totalizando $2N$ configurações. Ver Figura~\ref{fig:ordem_t8_completo}.
\end{itemize}
%
\begin{figure}[htbp]
    \centering
    % -------------------------------------------------------------
    % LINHA 1: GRAFOS DESCONECTADOS / TOCANTES
    % -------------------------------------------------------------
    \begin{subfigure}[b]{0.46\textwidth}
        \centering
        \begin{tikzpicture}[scale=0.7]
            % Rede de fundo
            \draw[gray!40, thin, dashed] (-0.3,-0.3) grid (3.3,2.3);
            \foreach \x in {0,...,3} {
                \foreach \y in {0,...,2} {
                    \fill[gray!50] (\x,\y) circle (2pt);
                }
            }
            % Quadrado 1
            \draw[teal!80!black, line width=1.5pt, line join=round] (0,0) rectangle (1,1);
            \foreach \x/\y in {0/0, 1/0, 1/1, 0/1} { \fill[teal!80!black] (\x,\y) circle (2.8pt); }
            % Quadrado 2 (disjunto)
            \draw[teal!80!black, line width=1.5pt, line join=round] (2,1) rectangle (3,2);
            \foreach \x/\y in {2/1, 3/1, 3/2, 2/2} { \fill[teal!80!black] (\x,\y) circle (2.8pt); }
        \end{tikzpicture}
        \caption{Dois quadrados disjuntos ($\frac{N^2}{2} - \frac{9N}{2}$).}
        \label{fig:t8_disjuntos}
    \end{subfigure}
    \hfill
    \begin{subfigure}[b]{0.46\textwidth}
        \centering
        \begin{tikzpicture}[scale=0.7]
            % Rede de fundo
            \draw[gray!40, thin, dashed] (-0.3,-0.3) grid (2.3,2.3);
            \foreach \x in {0,...,2} {
                \foreach \y in {0,...,2} {
                    \fill[gray!50] (\x,\y) circle (2pt);
                }
            }
            % Quadrados tocantes no vértice (1,1)
            \draw[violet!80!black, line width=1.5pt, line join=round] (0,0) rectangle (1,1);
            \draw[violet!80!black, line width=1.5pt, line join=round] (1,1) rectangle (2,2);
            \foreach \x/\y in {0/0, 1/0, 0/1, 1/1, 2/1, 2/2, 1/2} { \fill[violet!80!black] (\x,\y) circle (2.8pt); }
            % Destaque no vértice compartilhado (grau 4)
            \draw[violet!80!black, fill=white, line width=1.2pt] (1,1) circle (3.5pt);
        \end{tikzpicture}
        \caption{Tocantes por $1$ vértice ($2N$).}
        \label{fig:t8_vertice}
    \end{subfigure}

    \vspace{12pt}

    % -------------------------------------------------------------
    % LINHA 2: POLÍGONOS SIMPLES (LAÇOS DE COMPRIMENTO 8)
    % -------------------------------------------------------------
    \begin{subfigure}[b]{0.32\textwidth}
        \centering
        \begin{tikzpicture}[scale=0.65]
            % Rede de fundo
            \draw[gray!40, thin, dashed] (-0.3,-0.3) grid (3.3,1.3);
            \foreach \x in {0,...,3} {
                \foreach \y in {0,...,1} {
                    \fill[gray!50] (\x,\y) circle (2pt);
                }
            }
            % Retângulo 1x3
            \draw[orange!90!black, line width=1.5pt, line join=round] (0,0) -- (3,0) -- (3,1) -- (0,1) -- cycle;
            \foreach \x in {0,...,3} {
                \fill[orange!90!black] (\x,0) circle (2.8pt);
                \fill[orange!90!black] (\x,1) circle (2.8pt);
            }
        \end{tikzpicture}
        \caption{Retângulo $1\times 3$ ($2N$).}
        \label{fig:t8_1x3}
    \end{subfigure}
    \hfill
    \begin{subfigure}[b]{0.31\textwidth}
        \centering
        \begin{tikzpicture}[scale=0.65]
            % Rede de fundo
            \draw[gray!40, thin, dashed] (-0.3,-0.3) grid (2.3,2.3);
            \foreach \x in {0,...,2} {
                \foreach \y in {0,...,2} {
                    \fill[gray!50] (\x,\y) circle (2pt);
                }
            }
            % Quadrado 2x2
            \draw[blue!80!black, line width=1.5pt, line join=round] (0,0) -- (2,0) -- (2,2) -- (0,2) -- cycle;
            \foreach \x/\y in {0/0, 1/0, 2/0, 2/1, 2/2, 1/2, 0/2, 0/1} {
                \fill[blue!80!black] (\x,\y) circle (2.8pt);
            }
        \end{tikzpicture}
        \caption{Quadrado $2\times 2$ ($1N$).}
        \label{fig:t8_2x2}
    \end{subfigure}
    \hfill
    \begin{subfigure}[b]{0.32\textwidth}
        \centering
        \begin{tikzpicture}[scale=0.65]
            % Rede de fundo
            \draw[gray!40, thin, dashed] (-0.3,-0.3) grid (2.3,2.3);
            \foreach \x in {0,...,2} {
                \foreach \y in {0,...,2} {
                    \fill[gray!50] (\x,\y) circle (2pt);
                }
            }
            % Polígono em L (3 plaquetas)
            \draw[magenta!80!black, line width=1.5pt, line join=round] (0,0) -- (2,0) -- (2,1) -- (1,1) -- (1,2) -- (0,2) -- cycle;
            \foreach \x/\y in {0/0, 1/0, 2/0, 2/1, 1/1, 1/2, 0/2, 0/1} {
                \fill[magenta!80!black] (\x,\y) circle (2.8pt);
            }
        \end{tikzpicture}
        \caption{Polígono em "L" ($4N$).}
        \label{fig:t8_L}
    \end{subfigure}

    \caption{Classificação dos subgrafos eulerianos de ordem $t^8$ ($8$ arestas): (a) pares disjuntos, (b) pares tocantes e (c, d, e) laços simples de 8 arestas ($2N + 1N + 4N = 7N$). Somando todas as contribuições, obtém-se o coeficiente total de $\frac{N^2}{2} + \frac{9N}{2}$.}
    \label{fig:ordem_t8_completo}
\end{figure}
%
Somando todas as contribuições,
%
\begin{equation}
    \frac{N^2}{2} - \frac{9N}{2} + 7N + 2N
    = \frac{N^2}{2} + \frac{9N}{2}.
\end{equation}
%
Portanto, a expansão da função de partição até $t^8$ é
%
\begin{equation}
    \boxed{
    Z = (\cosh K)^{2N} 2^N
    \left[
    1 + N t^4 + 2N t^6
    + \left(\frac{N^2}{2} + \frac{9N}{2}\right) t^8
    + \mathcal O(t^{10})
    \right].
    }
    \label{eq:Z_expansion}
\end{equation}
% ------------------------------------------------------------
\subsection*{(b) Expansão da energia livre por spin}
% ------------------------------------------------------------
A energia livre por spin é
%
\begin{equation}
    g = \lim_{N\to\infty} \left(-\frac{1}{\beta N} \ln Z\right).
\end{equation}
%
Substituindo \eqref{eq:Z_expansion},
%
\begin{align}
    -\beta g
    &= \lim_{N\to\infty} \frac{1}{N} \ln Z \nonumber\\
    &= 2 \ln \cosh K + \ln 2
    + \lim_{N\to\infty} \frac{1}{N} \ln S,
\end{align}
%
onde
%
\begin{equation}
    S = 1 + N t^4 + 2N t^6
    + \left(\frac{N^2}{2} + \frac{9N}{2}\right) t^8
    + \mathcal O(t^{10}).
\end{equation}
%
Para expandir $\ln S$, utilizamos a série de Taylor genérica da função $\ln(1+x)$ em torno de $x = 0$:
\begin{equation}
    \ln(1 + x) = \sum_{k=1}^{\infty} \frac{(-1)^{k-1}}{k} x^k = x - \frac{1}{2} x^2 + \frac{1}{3} x^3 - \frac{1}{4} x^4 + \mathcal{O}(x^5),
    \label{eq:ln_taylor_generic}
\end{equation}
a qual é válida para $|x| < 1$. No nosso caso, escrevemos a soma sobre os grafos fechados $S$ na forma $S = 1 + x$, identificando a variável $x$ como a contribuição de todos os grafos não triviais:
%
\begin{equation}
    x \equiv N t^4 + 2N t^6 + \left(\frac{N^2}{2} + \frac{9N}{2}\right) t^8 + \mathcal{O}(t^{10}).
\end{equation}
%
Como queremos a expansão de $\ln S$ truncada até a ordem $t^8$, calculamos as potências de $x$ necessárias mantendo apenas os termos de ordem menor ou igual a $t^8$:
%
\begin{itemize}
    \item \textbf{Termo de ordem $x^1$:}
    %
    \begin{equation}
        x = N t^4 + 2N t^6 + \left(\frac{N^2}{2} + \frac{9N}{2}\right) t^8 + \mathcal{O}(t^{10}).
    \end{equation}
    %
    \item \textbf{Termo de ordem $x^2$:}
    %
    \begin{align}
        x^2 &= \left[ N t^4 + 2N t^6 + \left(\frac{N^2}{2} + \frac{9N}{2}\right) t^8 + \mathcal{O}(t^{10}) \right]^2 \nonumber\\
            &= \left(N t^4\right)^2 + 2(N t^4)(2N t^6) + \mathcal{O}(t^{10}) \nonumber\\
            &= N^2 t^8 + \mathcal{O}(t^{10}).
    \end{align}
    %
    \item \textbf{Termos de ordem $x^3, x^4, \dots$:}
    Como a menor potência em $x$ é $t^4$, temos $x^3 \sim \mathcal{O}(t^{12})$, de modo que todas as potências $x^k$ com $k \ge 3$ contribuem apenas com ordens superiores a $t^8$. 
\end{itemize}
%
Substituindo as potências $x$ e $x^2$ na expansão genérica \eqref{eq:ln_taylor_generic}:
%
\begin{align}
    \ln S &= \ln(1 + x) = x - \frac{1}{2}x^2 + \mathcal{O}(t^{10}) \nonumber\\[6pt]
          &= \left[ N t^4 + 2N t^6 + \left(\frac{N^2}{2} + \frac{9N}{2}\right) t^8 \right] - \frac{1}{2} \left[ N^2 t^8 \right] + \mathcal{O}(t^{10}) \nonumber\\[6pt]
          &= N t^4 + 2N t^6 + \left( \frac{N^2}{2} - \frac{N^2}{2} + \frac{9N}{2} \right) t^8 + \mathcal{O}(t^{10}) \nonumber\\[6pt]
          &= N t^4 + 2N t^6 + \frac{9N}{2} t^8 + \mathcal{O}(t^{10}).
\end{align}
%
O cancelamento exato do termo proporcional a $N^2$ no coeficiente de $t^8$ garante que a energia livre seja uma grandeza puramente extensiva (proporcional a $N$), permitindo um limite termodinâmico bem-definido e finito:
%
\begin{equation}
    \lim_{N\to\infty} \frac{1}{N} \ln S = t^4 + 2 t^6 + \frac{9}{2} t^8 + \mathcal{O}(t^{10}).
\end{equation}
%
Portanto,
%
\begin{equation}
    \boxed{
    -\beta g
    = 2 \ln \cosh K + \ln 2
    + t^4 + 2 t^6 + \frac{9}{2} t^8 + \mathcal O(t^{10}),
    }
    \label{eq:free_energy_expansion}
\end{equation}
%
que é exatamente o resultado fornecido no enunciado. Isso verifica os coeficientes da expansão de altas temperaturas da energia livre por spin.
% ============================================================
\section*{Problema 5 -- Método da Razão}
% ============================================================
A série fornecida para a função $f(x)$ é
%
\begin{equation}
    f(x)
    =
    4+16x+37x^2+67x^3+106x^4+154x^5
    +211x^6+277x^7+352x^8+436x^9+529x^{10}.
    \label{eq:f_series_p5}
\end{equation}
%
Supõe-se que, nas vizinhanças do ponto singular $x_c$, a função
apresente o comportamento assintótico
%
\begin{equation}
    f(x)\sim A(x_c-x)^{-\gamma}.
    \label{eq:singular_form_p5}
\end{equation}
%
O objetivo é utilizar o método da razão para obter estimativas para
$x_c$ e $\gamma$ e comparar os resultados com a solução analítica da
série. Consideremos inicialmente a expansão em série de potências
%
\begin{equation}
    (a-x)^b
    =
    \sum_{n=0}^{\infty} A_n x^n.
    \label{eq:binomial_p5}
\end{equation}
%
Escrevendo
%
\begin{equation}
    (a-x)^b
    =
    a^b
    \left(1-\frac{x}{a}\right)^b,
\end{equation}
%
e utilizando o teorema binomial generalizado,
%
\begin{equation}
    (1-z)^b
    =
    \sum_{n=0}^{\infty}
    (-1)^n
    \binom{b}{n}z^n,
\end{equation}
%
obtemos
%
\begin{equation}
    A_n
    =
    \frac{(-1)^n}{n!}
    b(b-1)(b-2)\cdots(b-n+1)
    a^{b-n}.
    \label{eq:An_p5}
\end{equation}
%
Consideremos agora a razão entre coeficientes consecutivos:
%
\begin{align}
    \frac{A_n}{A_{n-1}}
    &=
    \frac{
    \dfrac{(-1)^n}{n!}
    b(b-1)\cdots(b-n+1)a^{b-n}
    }{
    \dfrac{(-1)^{n-1}}{(n-1)!}
    b(b-1)\cdots(b-n+2)a^{b-n+1}
    }
    \nonumber\\
    &=
    -\frac{b-n+1}{an}
    \nonumber\\
    &=
    \frac{n-b-1}{an}
    \nonumber\\
    &=
    \frac{1}{a}
    \left[
        1-\frac{b+1}{n}
    \right].
\end{align}
%
Portanto,
%
\begin{equation}
    \frac{A_n}{A_{n-1}}
    =
    \frac{1}{a}
    \left[
        1-\frac{b+1}{n}
    \right].
    \label{eq:ratio_ab_p5}
\end{equation}
%
Para a singularidade considerada no problema,
%
\begin{equation}
    f(x)\sim A(x_c-x)^{-\gamma},
\end{equation}
%
fazemos a identificação
%
\begin{equation}
    a=x_c,
    \qquad
    b=-\gamma.
\end{equation}
%
Substituindo essas relações em \eqref{eq:ratio_ab_p5}, temos
%
\begin{equation}
    \frac{A_n}{A_{n-1}}
    =
    \frac{1}{x_c}
    \left[
        1+\frac{\gamma-1}{n}
    \right].
    \label{eq:ratio_final_p5}
\end{equation}
%
Assim, no limite assintótico $n\to\infty$,
%
\begin{equation}
    \frac{A_n}{A_{n-1}}
    =
    \frac{1}{x_c}
    +
    \frac{\gamma-1}{x_c}\frac{1}{n}
    +
    \mathcal O\left(\frac{1}{n^2}\right).
    \label{eq:ratio_asymptotic_p5}
\end{equation}
%
Essa expressão mostra que um gráfico de
%
\begin{equation}
    R_n\equiv\frac{A_n}{A_{n-1}}
\end{equation}
%
em função de $1/n$ deve se aproximar assintoticamente de uma reta. Seu
intercepto fornece $1/x_c$, enquanto sua inclinação fornece
$(\gamma-1)/x_c$. A partir da série \eqref{eq:f_series_p5}, identificamos
%
\begin{align}
    &A_0=4,\qquad
    A_1=16,\qquad
    A_2=37,\qquad
    A_3=67,\qquad
    A_4=106,
    \nonumber\\
    &A_5=154,\qquad
    A_6=211,\qquad
    A_7=277,\qquad
    A_8=352,
    \nonumber\\
    &A_9=436,\qquad
    A_{10}=529.
\end{align}
%
Uma característica importante da sequência é que suas diferenças
sucessivas são particularmente simples:
%
\begin{align}
    A_1-A_0 &=16-4=12,\\
    A_2-A_1 &=37-16=21,\\
    A_3-A_2 &=67-37=30,\\
    A_4-A_3 &=106-67=39.
\end{align}
%
Continuando,
%
\begin{equation}
    A_n-A_{n-1}=9n+3.
    \label{eq:difference_An_p5}
\end{equation}
%
Somando telescopicamente de $k=1$ até $n$,
%
\begin{align}
    A_n-A_0
    &=
    \sum_{k=1}^{n}(A_k-A_{k-1})
    \nonumber\\
    &=
    \sum_{k=1}^{n}(9k+3)
    \nonumber\\
    &=
    9\sum_{k=1}^{n}k
    +
    3\sum_{k=1}^{n}1.
\end{align}
%
Utilizando
%
\begin{equation}
    \sum_{k=1}^{n}k=\frac{n(n+1)}{2},
\end{equation}
%
obtemos
%
\begin{align}
    A_n-4
    &=
    9\frac{n(n+1)}{2}+3n
    \nonumber\\
    &=
    \frac{9n^2+9n+6n}{2}.
\end{align}
%
Logo,
%
\begin{align}
    A_n
    &=
    4+\frac{9n^2+15n}{2}
    \nonumber\\
    &=
    \frac{9n^2+15n+8}{2}.
\end{align}
%
Portanto, temos a expressão fechada
%
\begin{equation}
    A_n=\frac{9n^2+15n+8}{2}.
    \label{eq:An_closed_p5}
\end{equation}
%
Como
%
\begin{equation}
    f(x)
    =
    \sum_{n=0}^{\infty}A_nx^n,
\end{equation}
%
substituímos \eqref{eq:An_closed_p5}:
%
\begin{equation}
    f(x)
    =
    \sum_{n=0}^{\infty}
    \left(
        \frac{9n^2+15n+8}{2}
    \right)x^n.
\end{equation}
%
Separando os termos,
%
\begin{equation}
    f(x)
    =
    \frac{9}{2}\sum_{n=0}^{\infty}n^2x^n
    +
    \frac{15}{2}\sum_{n=0}^{\infty}nx^n
    +
    4\sum_{n=0}^{\infty}x^n.
    \label{eq:f_split_p5}
\end{equation}
%
Para $|x|<1$, temos as identidades
%
\begin{equation}
    \sum_{n=0}^{\infty}x^n
    =
    \frac{1}{1-x},
    \label{eq:sum0_p5}
\end{equation}
%
\begin{equation}
    \sum_{n=0}^{\infty}nx^n
    =
    \frac{x}{(1-x)^2},
    \label{eq:sum1_p5}
\end{equation}
%
e
%
\begin{equation}
    \sum_{n=0}^{\infty}n^2x^n
    =
    \frac{x(1+x)}{(1-x)^3}.
    \label{eq:sum2_p5}
\end{equation}
%
Substituindo essas expressões em \eqref{eq:f_split_p5},
%
\begin{equation}
    f(x)
    =
    \frac{9}{2}
    \frac{x(1+x)}{(1-x)^3}
    +
    \frac{15}{2}
    \frac{x}{(1-x)^2}
    +
    \frac{4}{1-x}.
\end{equation}
%
Colocando os três termos sobre o denominador comum $2(1-x)^3$,
%
\begin{align}
    f(x)
    &=
    \frac{
        9x(1+x)
        +
        15x(1-x)
        +
        8(1-x)^2
    }
    {2(1-x)^3}
    \nonumber\\
    &=
    \frac{
        9x+9x^2
        +
        15x-15x^2
        +
        8-16x+8x^2
    }
    {2(1-x)^3}
    \nonumber\\
    &=
    \frac{
        2x^2+8x+8
    }
    {2(1-x)^3}
    \nonumber\\
    &=
    \frac{
        x^2+4x+4
    }
    {(1-x)^3}.
\end{align}
%
Como
%
\begin{equation}
    x^2+4x+4=(x+2)^2,
\end{equation}
%
a função geradora exata é
%
\begin{equation}
    f(x)=\frac{(x+2)^2}{(1-x)^3}.    \label{eq:exact_generating_p5}
\end{equation}
%
A função \eqref{eq:exact_generating_p5} possui uma singularidade em
%
\begin{equation}
    x_c=1.
\end{equation}
%
Além disso, nas vizinhanças de $x=1$,
%
\begin{equation}
    (x+2)^2\longrightarrow9.
\end{equation}
%
Assim,
%
\begin{equation}
    f(x)
    \sim
    9(1-x)^{-3},
    \qquad
    x\to1.
\end{equation}
%
Comparando com a forma assumida
%
\begin{equation}
    f(x)\sim A(x_c-x)^{-\gamma},
\end{equation}
%
obtemos exatamente
%
\begin{equation}
    \boxed{
    x_c=1,
    \qquad
    \gamma=3,
    \qquad
    A=9.
    }
    \label{eq:exact_results_p5}
\end{equation}
%
Portanto, neste problema, os valores exatos de $x_c$ e $\gamma$ podem
ser determinados diretamente porque conhecemos a função geradora
fechada. A partir da expressão fechada para $A_n$, a razão consecutiva é
%
\begin{equation}
    R_n
    =
    \frac{A_n}{A_{n-1}}
    =
    \frac{
        9n^2+15n+8
    }{
        9(n-1)^2+15(n-1)+8
    }.
\end{equation}
%
Expandindo o denominador,
%
\begin{equation}
    9(n-1)^2+15(n-1)+8
    =
    9n^2-3n+2,
\end{equation}
%
e, portanto,
%
\begin{equation}
    R_n
    =
    \frac{9n^2+15n+8}
         {9n^2-3n+2}.
    \label{eq:R_exact_p5}
\end{equation}
%
Dividindo numerador e denominador por $9n^2$,
%
\begin{equation}
    R_n
    =
    \frac{
        1+\dfrac{5}{3n}+\dfrac{8}{9n^2}
    }{
        1-\dfrac{1}{3n}+\dfrac{2}{9n^2}
    }.
\end{equation}
%
Expandindo em potências de $1/n$, encontramos
%
\begin{equation}
    R_n
    =
    1
    +
    \frac{2}{n}
    +
    \frac{4}{3n^2}
    -
    \frac{8}{27n^4}
    +
    \mathcal O\left(\frac{1}{n^5}\right).    \label{eq:R_exact_asympt_p5}
\end{equation}
%
Em particular, mantendo apenas os termos dominantes,
%
\begin{equation}
    R_n
    =
    1+\frac{2}{n}
    +
    \mathcal O\left(\frac{1}{n^2}\right).
    \label{eq:R_dominant_p5}
\end{equation}
%
Comparando com a forma assintótica geral
%
\begin{equation}
    R_n
    =
    \frac{1}{x_c}
    +
    \frac{\gamma-1}{x_c}\frac{1}{n}
    +
    \mathcal O\left(\frac{1}{n^2}\right),
\end{equation}
%
segue que
%
\begin{align}
    \frac{1}{x_c}&=1,
    \\
    \frac{\gamma-1}{x_c}&=2.
\end{align}
%
Consequentemente,
%
\begin{equation}
    \boxed{
    x_c=1,
    \qquad
    \gamma=3,
    }
\end{equation}
%
em perfeito acordo com o resultado obtido diretamente da função
geradora. Vamos agora aplicar o método da razão apenas aos dez primeiros termos da série fornecida. Definimos
%
\begin{equation}
    R_n\equiv\frac{A_n}{A_{n-1}},
    \qquad
    n=1,\ldots,10.
\end{equation}
%
Os valores são apresentados na Tabela~\ref{tab:ratio_p5}.
%
\begin{table}[htbp]
    \centering
    \caption{
        Coeficientes $A_n$, variável $1/n$ e razões consecutivas
        $R_n=A_n/A_{n-1}$.
    }
    \label{tab:ratio_p5}
    \begin{tabular}{cccc}
        \toprule
        $n$ & $A_n$ & $1/n$ & $R_n$ \\
        \midrule
        1  & 16  & 1.0000 & 4.0000 \\
        2  & 37  & 0.5000 & 2.3125 \\
        3  & 67  & 0.3333 & 1.8108 \\
        4  & 106 & 0.2500 & 1.5821 \\
        5  & 154 & 0.2000 & 1.4528 \\
        6  & 211 & 0.1667 & 1.3701 \\
        7  & 277 & 0.1429 & 1.3128 \\
        8  & 352 & 0.1250 & 1.2708 \\
        9  & 436 & 0.1111 & 1.2386 \\
        10 & 529 & 0.1000 & 1.2133 \\
        \bottomrule
    \end{tabular}
\end{table}
%
Consideremos inicialmente a forma assintótica linear
%
\begin{equation}
    R_n
    \simeq
    c_0+c_1\frac{1}{n}.
    \label{eq:fit_linear_p5}
\end{equation}
%
Um ajuste linear utilizando todos os dez pontos, $n=1,\ldots,10$,
fornece
%
\begin{equation}
    \boxed{
    R_n
    \simeq
    0.8525
    +
    3.0861\frac{1}{n}.
    }
\end{equation}
%
Como o intercepto é aproximadamente $1/x_c$,
%
\begin{equation}
    x_c^{\rm (lin)}
    =
    \frac{1}{0.8525}
    \simeq
    1.1731.
\end{equation}
%
Da inclinação,
%
\begin{equation}
    \gamma^{\rm (lin)}
    =
    1+\frac{3.0861}{0.8525}
    \simeq
    4.6202.
\end{equation}
%
Portanto,
%
\begin{equation}
    \boxed{
    x_c^{\rm (lin)}=1.1731,
    \qquad
    \gamma^{\rm (lin)}=4.6202.
    }
    \label{eq:linear_results_p5}
\end{equation}
%
Esses valores diferem dos resultados exatos. A razão é que a fórmula
linear em $1/n$ é apenas uma aproximação assintótica, enquanto os
primeiros pontos, especialmente aqueles com $n$ pequeno, ainda possuem
correções importantes de ordem $1/n^2$ e ordens superiores. Para levar explicitamente em consideração a primeira correção
subdominante, consideramos
%
\begin{equation}
    R_n
    =
    c_0
    +
    c_1\frac{1}{n}
    +
    c_2\frac{1}{n^2}.
    \label{eq:fit_quad_p5}
\end{equation}
%
O ajuste aos dez pontos fornece
%
\begin{equation}
    \boxed{
    R_n
    \simeq
    0.9708
    +
    2.2904\frac{1}{n}
    +
    0.7404\frac{1}{n^2}.
    }
\end{equation}
%
Nesse caso,
%
\begin{equation}
    x_c^{\rm (quad)}
    =
    \frac{1}{0.9708}
    \simeq
    1.0301,
\end{equation}
%
e
%
\begin{equation}
    \gamma^{\rm (quad)}
    =
    1+\frac{2.2904}{0.9708}
    \simeq
    3.3593.
\end{equation}
%
Logo,
%
\begin{equation}
    \boxed{
    x_c^{\rm (quad)}=1.0301,
    \qquad
    \gamma^{\rm (quad)}=3.3593.
    }
    \label{eq:quad_results_p5}
\end{equation}
%
O ajuste quadrático produz estimativas consideravelmente mais próximas
dos valores exatos. Isso mostra que as correções de ordem $1/n^2$
desempenham um papel importante quando a série possui apenas alguns
termos. A análise anterior pode ser visualizada no gráfico de Domb--Sykes, que representa $R_n=\frac{A_n}{A_{n-1}}$ em função de $1/n$. A Fig.~\ref{fig:domb_sykes_p5} mostra os dados da série, os ajustes linear e quadrático e o comportamento assintótico exato
%
\begin{equation}
    R_n^{\rm ass}
    =
    1+\frac{2}{n}.
\end{equation}
%
\begin{figure}[htbp]
    \centering
    \includegraphics[width=0.88\linewidth]{Figures/statistical_mechanics/domb_sykes_problema5.pdf}
    \caption{
        Gráfico de Domb--Sykes para a série do Problema~5,
        mostrando as razões consecutivas
        $R_n=A_n/A_{n-1}$ em função de $1/n$.
        A linha tracejada (laranja) representa o ajuste linear
        $R_n=0.8525+3.0861/n$, que fornece
        $x_c^{\rm(lin)}=1.1731$ e $\gamma^{\rm(lin)}=4.6202$.
        A linha contínua (verde) corresponde ao ajuste quadrático
        $R_n=0.9708+2.2904/n+0.7404/n^2$, resultando em
        $x_c^{\rm(quad)}=1.0301$ e $\gamma^{\rm(quad)}=3.3593$.
        A linha pontilhada (vermelha) mostra o comportamento assintótico exato
        $R_n=1+2/n$, correspondente a $x_c=1$ e $\gamma=3$.
    }
    \label{fig:domb_sykes_p5}
\end{figure}
%
A figura evidencia que o ajuste linear aos pontos de baixa ordem ainda
apresenta um desvio significativo do comportamento assintótico. A
inclusão do termo em $1/n^2$ produz uma descrição mais próxima dos
dados e leva as estimativas para mais perto dos valores exatos. É importante destacar que a linha $1+2/n$ representa o comportamento
assintótico da razão, e não os valores exatos de $R_n$ para $n$ finito. A expressão exata contém as correções subdominantes mostradas em \eqref{eq:R_exact_asympt_p5}. Os resultados podem ser resumidos na Tabela~\ref{tab:fit_results_p5}.
%
\begin{table}[htbp]
    \centering
    \caption{
        Comparação entre as estimativas obtidas pelo método da razão e
        os valores exatos determinados pela função geradora.
    }
    \label{tab:fit_results_p5}
    \begin{tabular}{lcc}
        \toprule
        Método & $x_c$ & $\gamma$ \\
        \midrule
        Ajuste linear
        & $1.1731$
        & $4.6202$
        \\
        Ajuste quadrático
        & $1.0301$
        & $3.3593$
        \\
        Resultado exato
        & $1.0000$
        & $3.0000$
        \\
        \bottomrule
    \end{tabular}
\end{table}
%
A comparação mostra claramente a diferença entre uma análise
assintótica aplicada a uma série curta e a determinação exata dos
parâmetros. O ajuste linear fornece uma aproximação relativamente pobre para os valores de $x_c$ e $\gamma$, pois os termos de baixa ordem ainda contêm contribuições importantes das correções em potências de $1/n$. Ao incluir um termo quadrático, as estimativas melhoram.
% ============================================================
\section*{Problema 6 -- Modelo autodual}
% ============================================================
Consideremos o modelo de Ising ferromagnético na rede quadrada, em campo externo nulo, com hamiltoniano
%
\begin{equation}
    \mathcal H
    =
    -J\sum_{\langle ij\rangle}\sigma_i\sigma_j,
    \qquad
    J>0,
    \qquad
    \sigma_i=\pm1.
    \label{eq:H_ising_2d_dual}
\end{equation}
%
Definimos, como usual,
%
\begin{equation}
    K\equiv\beta J=\frac{J}{k_BT}.
    \label{eq:K_def}
\end{equation}
%
O objetivo é mostrar a correspondência um-a-um entre as expansões de altas e
baixas temperaturas e obter a temperatura crítica pelo argumento de
autodualidade. A função de partição é
%
\begin{equation}
    Z(K)
    =
    \sum_{\{\sigma\}}
    \exp\left(
        K\sum_{\langle ij\rangle}
        \sigma_i\sigma_j
    \right).
    \label{eq:Z_high_start}
\end{equation}
%
Para cada ligação, usamos a identidade, já derivada no problema 4,
%
\begin{equation}
    e^{K\sigma_i\sigma_j}
    =
    \cosh K
    \left(
        1+\sigma_i\sigma_j\tanh K
    \right).
    \label{eq:high_temp_identity}
\end{equation}
%
Definindo
%
\begin{equation}
    t\equiv\tanh K,
    \label{eq:t_def}
\end{equation}
%
podemos escrever
%
\begin{equation}
    Z(K)
    =
    (\cosh K)^E
    \sum_{\{\sigma\}}
    \prod_{\langle ij\rangle}
    \left(
        1+t\,\sigma_i\sigma_j
    \right),
    \label{eq:high_temp_expansion}
\end{equation}
%
onde \(E\) é o número total de ligações da rede. Ao expandir o produto, aparecem produtos de spins associados a subconjuntos de ligações. A soma sobre cada variável de spin \(\sigma_i\) se anula se \(\sigma_i\) aparecer com potência ímpar, pois
%
\begin{equation}
    \sum_{\sigma_i=\pm1}\sigma_i=0,
    \qquad
    \sum_{\sigma_i=\pm1}1=2.
\end{equation}
%
Portanto, somente sobrevivem os subconjuntos de ligações nos quais cada sítio tem número par de ligações selecionadas. Na linguagem gráfica, esses subconjuntos formam \textbf{laços fechados} (subgrafos de Euler) na rede original. Assim,
%
\begin{equation}
    Z(K)
    =
    2^N
    (\cosh K)^E
    \sum_{\mathcal G}
    (\tanh K)^{\ell(\mathcal G)},
    \label{eq:Z_high_loops}
\end{equation}
%
onde \(N\) é o número de sítios, \(\mathcal G\) representa uma configuração de laços fechados e \(\ell(\mathcal G)\) é o número de ligações dessa configuração. Para a rede quadrada $E=2N$. Consideremos agora o modelo dual com acoplamento \(K^*\), a baixas temperaturas. No estado fundamental, todos os spins estão alinhados e cada ligação contribui com \(-J\). Existem dois estados fundamentais globalmente degenerados,
%
\begin{equation}
    \sigma_i=+1
    \qquad\text{ou}\qquad
    \sigma_i=-1.
\end{equation}
%
Uma excitação é produzida quando alguns spins são invertidos. Cada ligação que separa spins alinhados de spins antiparalelos passa de \(-J\) para \(+J\), resultando em um custo de energia
%
\begin{equation}
    \Delta E=2J.
\end{equation}
%
Portanto, uma configuração com \(\ell\) ligações de domínio tem fator de Boltzmann
%
\begin{equation}
    e^{-2K^*\ell}.
\end{equation}
%
As paredes de domínio formam caminhos fechados na \textbf{rede dual}. Assim, para o modelo dual,
%
\begin{equation}
    Z^*(K^*)
    =
    2\,e^{K^* E}
    \sum_{\mathcal G^*}
    e^{-2K^*\ell(\mathcal G^*)},
    \label{eq:low_temp_expansion}
\end{equation}
%
onde \(\mathcal G^*\) representa uma configuração de paredes de domínio na rede dual. O fator \(2\) corresponde aos dois estados fundamentais relacionados pela simetria global \(\sigma_i\to-\sigma_i\). Na rede quadrada, a rede dual é também uma rede quadrada. Existe uma correspondência um-a-um entre os laços fechados da expansão de altas temperaturas da rede original e as paredes de domínio da expansão de baixas temperaturas da rede dual. Os mesmos objetos geométricos aparecem nas duas expansões, mas com pesos diferentes:
%
\begin{equation}
    (\tanh K)^{\ell}
    \qquad\longleftrightarrow\qquad
    e^{-2K^*\ell}.
\end{equation}
%
Para que os termos correspondentes tenham o mesmo peso, impomos
%
\begin{equation}
    \tanh K=e^{-2K^*}.
    \label{eq:duality_relation_1}
\end{equation}
%
Como a transformação deve ser invertível, também vale
%
\begin{equation}
    \tanh K^*=e^{-2K}.
    \label{eq:duality_relation_2}
\end{equation}
%
Essas relações definem a transformação de dualidade \(K\leftrightarrow K^*\). Uma forma simétrica útil é
%
\begin{equation}
    \sinh(2K)\,\sinh(2K^*)=1.
    \label{eq:sinh_duality}
\end{equation}
%
Essa é a transformação de dualidade de Kramers--Wannier. Na rede quadrada, a rede dual tem a mesma estrutura da rede original; o
modelo é \textbf{autodual}. Se houver uma única temperatura crítica bem definida, ela deve ser invariante sob a dualidade. No ponto crítico,
%
\begin{equation}
    K^*=K=K_c.
    \label{eq:self_dual_condition}
\end{equation}
%
Substituindo em \eqref{eq:duality_relation_1},
%
\begin{equation}
    e^{-2K_c}=\tanh K_c.
    \label{eq:critical_duality}
\end{equation}
%
Como \(K_c=J/(k_BT_c)\), segue
%
\begin{equation}
    \boxed{
    \exp\left(
        -\frac{2J}{k_BT_c}
    \right)
    =
    \tanh\left(
        \frac{J}{k_BT_c}
    \right),
    }
    \label{eq:Tc_equation}
\end{equation}
%
que é exatamente a relação pedida. Definimos
%
\begin{equation}
    x\equiv e^{2K_c}.
\end{equation}
%
Então
%
\begin{equation}
    e^{-2K_c}=\frac1x,
    \qquad
    \tanh K_c
    =
    \frac{\sinh K_c}{\cosh K_c}
    =
    \frac{e^{K_c}-e^{-K_c}}{e^{K_c}+e^{-K_c}}
    =
    \frac{e^{2K_c}-1}{e^{2K_c}+1}
    =
    \frac{x-1}{x+1}.
\end{equation}
%
A condição crítica torna-se
%
\begin{equation}
    \frac1x
    =
    \frac{x-1}{x+1}.
\end{equation}
%
Multiplicando por \(x(x+1)\),
%
\begin{equation}
    x+1=x(x-1)
    \quad\Longrightarrow\quad
    x^2-2x-1=0.
\end{equation}
%
As soluções são \(x=1\pm\sqrt2\). Como \(x=e^{2K_c}>0\), escolhemos
%
\begin{equation}
    x=1+\sqrt2.
\end{equation}
%
Portanto,
%
\begin{equation}
    e^{2K_c}=1+\sqrt2,
\end{equation}
%
e finalmente
%
\begin{equation}
    K_c
    =
    \frac12\ln(1+\sqrt2).
    \label{eq:Kc_exact}
\end{equation}
%
Como \(K_c=J/(k_BT_c)\),
%
\begin{equation}
    \frac{J}{k_BT_c}
    =
    \frac12\ln(1+\sqrt2).
    \label{eq:critical_temperature_dimensionless}
\end{equation}
%
Equivalentemente,
%
\begin{equation}
    \frac{k_BT_c}{J}
    =
    \frac{2}{\ln(1+\sqrt2)}
    \simeq 2.269185.
    \label{eq:Tc_exact_dimensionless}
\end{equation}
%
Numericamente,
%
\begin{equation}
    K_c\simeq0.440687.
\end{equation}
%
A solução exata de Onsager para a rede quadrada em campo nulo fornece
%
\begin{equation}
    \sinh(2K_c)=1.
    \label{eq:onsager_condition}
\end{equation}
%
Tomando a função inversa,
%
\begin{equation}
    2K_c=\operatorname{arsinh}(1).
\end{equation}
%
Como
%
\begin{equation}
    \operatorname{arsinh}(1)
    =
    \ln(1+\sqrt2),
\end{equation}
%
obtemos
%
\begin{equation}
    K_c=\frac12\ln(1+\sqrt2),
\end{equation}
%
exatamente o mesmo resultado obtido pelo argumento de dualidade. Portanto, para o modelo de Ising bidimensional na rede quadrada, o ponto autodual coincide com a temperatura crítica exata:
%
\begin{equation}
    \boxed{
    \left(\frac{k_BT_c}{J}\right)_{\rm dualidade}
    =
    \left(\frac{k_BT_c}{J}\right)_{\rm Onsager}
    =
    \frac{2}{\ln(1+\sqrt2)}
    \simeq2.269185.
    }
\end{equation}
%
Na aproximação de campo médio, cada spin sente um campo efetivo produzido
pelos seus \(z\) primeiros vizinhos. Para a rede quadrada, \(z=4\). A equação de autoconsistência para a magnetização é
%
\begin{equation}
    m
    =
    \tanh(\beta zJm).
    \label{eq:mf_self_consistency}
\end{equation}
%
Próximo da transição, \(m\) é pequeno e
%
\begin{equation}
    \tanh x=x-\frac{x^3}{3}+\cdots.
\end{equation}
%
Mantendo apenas o termo linear,
%
\begin{equation}
    m\simeq\beta zJm.
\end{equation}
%
Para surgir uma solução não trivial \(m\neq0\), devemos ter
\begin{equation}
    1=\beta_c zJ.
\end{equation}
%
Logo,
%
\begin{equation}
    K_c^{\rm MF}
    =
    \beta_cJ
    =
    \frac1z.
\end{equation}
%
Para a rede quadrada,
%
\begin{equation}
    \boxed{
    K_c^{\rm MF}=\frac14.
    }
\end{equation}
%
Consequentemente,
%
\begin{equation}
    \left(\frac{k_BT_c}{J}\right)_{\rm MF}
    =4.
    \label{eq:Tc_MF}
\end{equation}
%
\begin{table}[H]
    \centering
    \caption{Temperatura crítica do modelo de Ising na rede quadrada em
    diferentes aproximações.}
    \label{tab:Tc_comparison}
    \begin{tabular}{c c c}
        \toprule
        Método & \(K_c=J/(k_BT_c)\) & \(k_BT_c/J\) \\
        \midrule
        Dualidade / Onsager
        & \(\dfrac12\ln(1+\sqrt2)\)
        & \(2.269185\) \\[0.3cm]

        Campo médio
        & \(\dfrac14\)
        & \(4\) \\
        \bottomrule
    \end{tabular}
\end{table}
%
A aproximação de campo médio fornece uma temperatura crítica diferente da
obtida pela solução exata. A origem dessa diferença está no fato de que a
aproximação de campo médio substitui as correlações entre os spins por um
campo efetivo médio, enquanto a solução exata bidimensional incorpora as flutuações e correlações espaciais do sistema. A expansão de altas temperaturas do modelo de Ising pode ser representada por configurações de laços fechados na rede original. A expansão de baixas temperaturas pode ser representada pelas paredes de domínio, que formam os mesmos tipos de laços na rede dual. 
% ============================================================
\end{document}
